\documentclass[12pt,a4paper]{ctexart}
\usepackage[top=2.5cm,bottom=2.5cm,left=3cm,right=3cm]{geometry}
\usepackage{graphicx}
\usepackage{fancyhdr}
\usepackage{titlesec}
\usepackage{enumitem}
\usepackage{listings}
\usepackage{hyperref}

% 去除默认目录名
\ctexset{contentsname={}}

% 页眉页脚
\pagestyle{fancy}
\fancyhf{}
\fancyhead[L]{\small 审计智鉴流式数据处理系统 V1.0}
\fancyhead[R]{\small \leftmark}
\fancyfoot[C]{\thepage}
\renewcommand{\headrulewidth}{0.4pt}
\renewcommand{\footrulewidth}{0pt}
\renewcommand{\sectionmark}[1]{\markboth{\thesection\quad #1}{}}

% 章节标题样式
\titleformat{\section}{\heiti\zihao{-2}\bfseries}{\thesection}{1em}{}
\titleformat{\subsection}{\heiti\zihao{3}\bfseries}{\thesubsection}{1em}{}
\titleformat{\subsubsection}{\heiti\zihao{4}\bfseries}{\thesubsubsection}{1em}{}

% 列表样式
\setlist[itemize]{leftmargin=2em,itemsep=0.3em}
\setlist[enumerate]{leftmargin=2em,itemsep=0.3em}

% 代码样式：纯黑白，无任何彩色
\lstset{
    language=Python,
    basicstyle=\ttfamily\small,
    frame=single,
    breaklines=true,
    showstringspaces=false,
    tabsize=4
}

% 截图占位框
\newcommand{\screenshotplaceholder}[2]{%
    \begin{center}
    \fbox{\parbox{0.85\textwidth}{%
        \centering
        \vspace{1.5cm}
        {\large \textbf{【截图位置：#1】}}\\[0.5em]
        {\small #2}
        \vspace{1.5cm}
    }}
    \end{center}
}

\begin{document}

% ==================== 封面 ====================
\begin{titlepage}
\centering
\vspace*{3cm}
{\Huge\heiti\bfseries 审计智鉴流式数据处理系统}\\[0.8cm]
{\LARGE\heiti 操作说明书}\\[0.5cm]
{\Large V1.0}\\[3cm]
\begin{tabular}{rl}
\textbf{软件全称：} & 审计智鉴流式数据处理系统 V1.0 \\[0.3em]
\textbf{软件简称：} & 审计智鉴系统 \\[0.3em]
\textbf{版本号：} & V1.0 \\[0.3em]
\textbf{编制日期：} & 2026年01月 \\
\end{tabular}
\vfill
\end{titlepage}

% ==================== 目录 ====================
\tableofcontents
\newpage

% ==================== 第1章 系统概述 ====================
\section{系统概述}

\subsection{系统用途与背景}

财务舞弊是资本市场长期存在的一类风险事件。康美药业、康得新、瑞幸咖啡等典型案例表明，传统人工审阅年报的方式在效率和深度方面均存在局限，大量风险信号隐藏在文本细节中，难以通过人工方式全面识别。

审计智鉴系统的开发目标在于构建一套智能化财务数据处理工具，通过程序辅助完成财报审阅工作，帮助审计人员快速定位可疑点，将主要精力集中于需要深度核查的领域。

系统核心解决两类问题：一是传统财务指标分析（如存贷双高、现金流背离等），二是管理层讨论与分析（MD\&A）文本中的语义风险识别。将二者结合，形成较为完整的舞弊风险评估体系。

\subsection{技术架构}

系统采用前后端分离架构，主要基于开发效率与部署灵活性进行技术选型。

后端采用 \textbf{FastAPI} 框架，基于 Python 构建异步 Web 服务，天然支持异步 IO，适用于调用外部 AI 接口的场景。后端承担全部核心业务逻辑，包括用户认证、检测分析、AI 问答、报告生成、会员计费等模块。

前端采用 \textbf{Streamlit} 框架。该框架面向数据类应用，开发效率较高，且对数据可视化支持良好，图表、表格、进度条等组件可直接调用。对于工具型应用而言，其 UI 自定义能力已可满足基本需求。

数据层采用 \textbf{MySQL} 数据库，通过 SQLAlchemy ORM 进行操作。缓存采用内存缓存方案，适用于当前用户规模。

AI 能力接入阿里云百炼（DashScope）通义千问模型，支持流式输出。选型依据为中文理解能力较强，且 API 接口采用 OpenAI 兼容格式，便于对接。

\subsection{系统核心能力}

系统具备以下核心能力：

\begin{enumerate}
    \item \textbf{双模融合检测}：传统财务指标占 40\% 权重，AI 文本分析占 60\% 权重，综合计算舞弊概率。
    \item \textbf{七维度文本风险分析}：从语义矛盾、风险披露、语调异常、文本-数据一致性、关联隐藏、信息密度、回避表述七个角度评估 MD\&A 文本。
    \item \textbf{SHAP 可解释性}：基于博弈论的 Shapley 值方法，量化各特征对最终判断的贡献度，满足审计报告可解释性要求。
    \item \textbf{流式 AI 问答}：采用 SSE（Server-Sent Events）技术，AI 回答内容逐字实时呈现。
    \item \textbf{多格式报告导出}：检测结果支持导出为 PDF、Word、Excel 三种格式。
\end{enumerate}

\newpage

% ==================== 第2章 运行环境 ====================
\section{运行环境说明}

\subsection{服务器环境}

系统部署于云服务器环境，操作系统为 Linux Ubuntu 22.04 LTS。该版本稳定性较好，社区支持成熟。

服务器基础配置要求为 2 核 CPU、4GB 内存。主要性能瓶颈在于外部 AI 接口的调用并发，而非系统本身计算压力。核心 XGBoost 模型推理速度为毫秒级。

\subsection{访问方式}

系统无专属域名，用户通过公网 IP 地址加端口号访问。后端 API 服务运行于 8000 端口，前端 Streamlit 界面运行于 8501 端口。该部署方式省去了域名注册与备案流程。

前后端分离架构下，前端页面加载后，所有数据交互均通过 RESTful API 完成。用户在前端上传文件、提交检测请求，前端将请求发送至后端，后端处理完成后返回结果，前端负责渲染展示。

\subsection{运行依赖}

系统运行需满足以下条件：

\begin{itemize}
    \item Python 3.9 或更高版本
    \item MySQL 8.0 数据库
    \item 阿里云 DashScope API 访问权限（用于 AI 分析）
    \item 充足磁盘空间用于存储上传文件及生成报告
\end{itemize}

数据库初始化在系统启动时自动完成，首次启动时 SQLAlchemy 需检查和同步表结构，耗时约数秒。

\newpage

% ==================== 第3章 系统功能模块 ====================
\section{系统功能模块}

本章对系统五个主要功能模块进行说明。各模块均围绕"数据接收—运算处理—结果输出"流程展开，前后端通过接口联动完成业务闭环。

\subsection{用户认证模块}

该模块负责用户注册、登录与权限管理。密码采用 PBKDF2-SHA256 算法哈希存储。登录成功后后端签发 JWT Token，有效期 7 天，前端将 Token 存储于 Cookie 中，后续请求自动携带。

用户分为免费版、专业版、企业版三个等级，不同等级在检测次数和 AI 问答次数上存在差异：

\begin{itemize}
    \item 免费版：每月 3 次检测额度，每天 5 次 AI 问答
    \item 专业版：检测不限次数，AI 问答每天 50 次
    \item 企业版：检测与 AI 问答均不限量
\end{itemize}

\subsection{舞弊检测模块}

该模块为系统核心。用户可通过三种方式提交检测：

\begin{enumerate}
    \item \textbf{文件上传}：上传 Excel 或 CSV 格式财务数据文件，系统自动解析年份及各项指标。可补充上传 MD\&A 文本文件（txt、docx、pdf 格式），用于文本风险分析。
    \item \textbf{内置案例库}：系统预置若干经典案例（如康美药业、贵州茅台等），用户可直接加载案例数据进行检测，便于快速体验功能。
    \item \textbf{手动录入}：在页面上直接填写各项财务指标，适用于数据量不大的场景。
\end{enumerate}

检测流程如下：后端收到请求后，先执行传统财务指标分析（存贷双高、现金流背离、存货异常等），再调用 AI 模型对 MD\&A 文本进行七维度分析，最后将两部分得分加权汇总，计算 0 至 100 的风险评分及舞弊概率。完整流程耗时约 3 至 5 秒，主要取决于外部 AI 接口响应速度。

\subsection{AI 问答模块}

该模块为流式输出智能问答功能。用户在输入框提问后，系统通过 SSE 流式返回 AI 回答，文字逐字呈现。

流式输出在技术实现上采用 Server-Sent Events。前端建立长连接，后端通过 httpx 的 stream 模式读取 AI 接口的流式响应，逐段推送给前端。相较于传统"等待全部内容生成后再返回"的方式，该方案用户体验更优。

系统预设五个类别的问题推荐：理论类、实践类、政策类、案例类、平台使用类。用户可在侧面栏直接点击推荐问题，无需手动输入。

\subsection{财务助手模块}

财务助手包含三个子功能：

\begin{itemize}
    \item \textbf{税务测算}：按最新个税法和企业所得税法计算应纳税额。个税测算支持月薪、年终奖、五险一金、专项附加扣除等，可对比年终奖"单独计税"与"合并计税"两种方式。企业税测算支持小规模纳税人和一般纳税人两种类型，区分不同行业增值税税率。
    \item \textbf{简易报税辅助}：用户上传增值税发票图片，系统调用多模态 AI 模型识别发票内容，提取金额、税率、税额等信息，自动生成报税数据汇总表。
    \item \textbf{财税咨询}：侧重税务政策解读与实操建议，同样采用流式输出。
\end{itemize}

\subsection{报告管理模块}

检测完成后，用户可生成正式检测报告。报告类型分基础版、专业版、企业版三种：

\begin{itemize}
    \item 基础版：检测结果摘要与风险标签
    \item 专业版：增加 SHAP 特征分析解读、整改建议、IPO 对标分析
    \item 企业版：包含完整技术细节与原始数据附录
\end{itemize}

报告支持导出为 PDF、Word、Excel 三种格式。PDF 采用预定义模板，排版较为正式；Word 便于二次编辑；Excel 将数据表格单独导出，便于进一步分析。

\newpage

% ==================== 第4章 系统操作流程 ====================
\section{系统操作流程}

本章通过实际操作流程说明系统使用方法。以下操作均在前端交互界面完成，后端通过 API 接口实时响应。

\subsection{首次使用流程}

首次打开系统，显示首页。页面上方为渐变色装饰条，中间为品牌 Logo 与导航菜单。首页主体为深色背景大图区，包含系统核心定位说明及动态统计数字（AI 识别准确率、已分析财报数量等），数字在页面滚动至可视区域时执行从 0 滚动至目标值的动画效果。

\screenshotplaceholder{img-01}{此处应插入系统主界面截图，需包含顶部导航栏、Hero 区域和动态统计数据栏}

新用户需先注册账号。点击右上角"登录 / 注册"按钮，弹出登录弹窗。弹窗分登录和注册两个标签页，注册时填写用户名、密码、邮箱（可选）、手机号（可选）即可。注册成功后自动登录，跳转至首页。

\screenshotplaceholder{img-02}{此处应插入登录/注册弹窗截图}

\subsection{执行舞弊检测}

登录后，点击导航栏"舞弊检测"进入检测页面。页面上方为企业名称和证券代码输入框，下方为年份区间设置，再下方为文件上传区域。

\screenshotplaceholder{img-03}{此处应插入舞弊检测-文件上传界面截图}

以文件上传方式为例，操作步骤如下：

\begin{enumerate}
    \item 填写企业名称（必填）和证券代码（可选）。
    \item 设置数据年份区间，系统默认 2020 至 2023 年。
    \item 上传财务数据文件，支持 xlsx、xls、csv 格式。若文件包含"年份"列，系统自动按年份分组解析；若无年份列，则将整个文件作为起始年份数据。
    \item （可选）上传 MD\&A 文本文件，支持 txt、docx、doc、pdf 格式，可多选。
    \item 点击"解析文件内容"按钮，系统解析文件并在页面显示预览结果，包括各年度提取的财务指标。
    \item 确认数据无误后，点击"开始批量检测"按钮，系统逐年度执行检测分析，页面显示进度条。
\end{enumerate}

检测完成后，页面自动跳转至结果展示页。结果页顶部为综合信息卡片，显示舞弊概率、风险评分、风险标签数量、置信度四项核心指标。下方左侧为舞弊概率仪表盘图表（环形图），右侧为风险标签统计。

\screenshotplaceholder{img-04}{此处应插入舞弊检测-检测结果截图，需包含风险仪表盘和风险标签}

再往下为 AI 风险特征雷达图，七个维度（语义矛盾度、风险披露完整性、异常乐观语调、文本-数据一致性、关联隐藏指数、信息密度异常、回避表述强度）得分一目了然。右侧为 SHAP 特征重要性排行，以进度条形式展示 Top 5 特征对模型判断的贡献度。

\subsection{使用 AI 问答（流式输出效果）}

点击导航栏"AI 问答"进入问答页面。首次进入时，页面中央显示欢迎卡片，介绍该功能为财务舞弊领域专业问答助手。

在底部输入框输入问题，如"康美药业舞弊案的关键识别点是什么"，按回车发送。页面显示用户问题气泡，AI 助手位置出现三个跳动圆点（thinking 动画），表示后端正在处理。

约一至两秒后，回答文字开始逐字出现，即 SSE 流式输出效果。文字并非一次性全部呈现，而是逐字符显示。该体验与传统网页"提交表单、等待加载、整页刷新"模式存在本质区别，更接近实时运行的应用程序。

\screenshotplaceholder{img-05}{此处应插入 AI 问答流式输出效果截图，需能看到文字逐字出现的状态。这是证明系统不是静态网页的关键截图，必须包含}

流式输出结束后，完整回答保留在对话记录中。用户可继续追问，系统支持多轮上下文对话。所有历史问答记录可在侧边栏"历史记录"中查看，亦可收藏常用问答。

\subsection{使用税务测算}

进入"财务助手"页面，点击"税务中心"标签。该页面分三个子标签：税务测算、简易报税、财税咨询。

以个税测算为例：填写税前月薪、年终奖金额、每月五险一金缴纳额、专项附加扣除额，点击"计算个税"按钮。系统立即计算年收入、年扣除额、适用税率、年个税、税后年收入五项指标。下方展开区域显示应纳税所得额详细计算过程，包括基本减除费用 6 万元、五险一金年总额、专项附加扣除年总额等逐项明细。

若填写了年终奖，系统自动对比"单独计税"与"合并计税"两种方式，推荐更省税方案并显示节省金额。

\screenshotplaceholder{img-06}{此处应插入税务测算结果截图}

\subsection{生成并导出报告}

检测完成后，在"我的检测"页面可查看所有历史检测记录。点击某条记录旁"查看详情"，进入详细结果页。页面上方有"导出报告"按钮，点击后选择报告类型（基础版/专业版/企业版）和导出格式（PDF/Word/Excel），系统在后端生成文件并提供下载链接。

\screenshotplaceholder{img-07}{此处应插入报告导出界面截图}

报告生成耗时约 5 至 10 秒，需将所有检测数据、图表、分析文本整合至一个文档中。PDF 版本采用预定义模板，带有页眉页脚和封面，可直接用于工作汇报。

\subsection{会员订阅}

免费版用户检测额度用完后，系统将提示升级会员。点击导航栏"会员"进入定价页面。页面展示三个版本会员方案：基础版（299元/年）、专业版（999元/年）、企业版（按需定价），各方案下列出包含功能与额度。

\screenshotplaceholder{img-08}{此处应插入会员中心定价页面截图}

用户选择方案后，点击"立即订阅"按钮，系统创建订单并跳转至支付页面。目前支持支付宝和微信支付两种渠道。支付成功后会员权益立即生效，无需重新登录。

\newpage

% ==================== 第5章 核心代码逻辑说明 ====================
\section{核心代码逻辑说明}

\subsection{前后端联动流程}

系统前后端交互遵循标准请求-响应模式，AI 问答模块采用长连接流式传输，本节对此进行单独说明。

普通接口（如用户登录、获取检测历史）的流程为：前端使用 Python requests 库发送 HTTP 请求，后端 FastAPI 处理请求、操作数据库、返回 JSON 数据。前后端数据格式统一采用 Pydantic Schema 进行校验。

AI 问答流式接口的差异在于：前端发送 POST 请求时增加 \texttt{stream=True} 参数。后端收到请求后，调用 httpx 异步流式客户端向阿里云 DashScope API 发起流式请求。DashScope 返回的响应为分段 SSE 数据（格式为 \texttt{data: \{...\}}），后端将数据片段原样转发给前端。前端通过 \texttt{response.iter\_lines()} 逐行读取，解析出文字内容，再使用 Streamlit 的 \texttt{st.write\_stream()} 方法实现逐字渲染。

该流程包含两个技术关键点：一是 httpx 的异步流式支持，使后端可在接收 AI 接口数据的同时转发给前端，无需等待全部内容；二是 Streamlit 的 \texttt{write\_stream} 方法，其内部做了增量渲染优化，仅更新变化部分，避免整段文字重绘导致闪烁。

\subsection{风险评分算法}

舞弊概率计算采用双模融合策略，公式如下：

\begin{center}
\fbox{\parbox{0.8\textwidth}{%
\centering
总风险评分 = 基础风险分（35）+ 传统财务风险评分（0--40）+ AI 文本风险评分（0--65）\\[0.3em]
舞弊概率 = 0.50 + (总风险评分 / 100) $\times$ 0.5
}}
\end{center}

传统财务风险评分主要考察三项经典指标：存贷双高（货币资金和短期借款同时处于高位）、现金流与利润背离（净利润为正但经营现金流为负）、存货占比异常（存货占总资产比例过高）。各指标达到阈值时累加对应分数，总分封顶 40 分。

AI 文本风险评分从 MD\&A 文本中提取七个维度特征得分，各维度具有不同权重。"文本-数据一致性"与"语义矛盾度"权重最高（2.0 倍），原因在于这两个维度对舞弊判断的区分度最强。所有维度加权后归一化至 65 分。

基础风险分设为 35 分而非 0 分，系出于审计场景的审慎考虑——财务数据本身存在不完全可信性，即使未出现明显异常，亦应保留一定基础警惕性。最终舞弊概率限制在 50\% 至 95\% 之间，最低不低于 50\%。

\subsection{SHAP 可解释性分析}

模型的可解释性对于审计类应用具有重要意义。仅告知用户"某企业存在 80\% 舞弊概率"是不够的，还需解释该结论的推导依据，即哪些因素推动了该判断。

系统采用 SHAP（SHapley Additive exPlanations）方法计算各特征对最终预测结果的贡献度。SHAP 值的核心思想源自博弈论：将每个特征视为一名"参与者"，计算其加入"联盟"时对预测结果的边际贡献。该方法具有良好的数学性质，即所有特征的 SHAP 值之和等于预测值与基准值之差，满足可加性解释。

实际实现中，由于采用 XGBoost 模型，可直接调用 \texttt{shap.TreeExplainer} 计算 SHAP 值，速度较快。计算结果包含正值（推动舞弊概率升高）与负值（推动舞弊概率降低），系统取绝对值最大的前 5 个特征展示给用户，并以进度条形式可视化。

\newpage

% ==================== 第6章 系统运行总结 ====================
\section{系统运行总结}

审计智鉴系统从立项至今，功能已较为完整。以下对设计决策进行回顾与总结。

采用 Streamlit 作为前端框架，主要基于开发效率的考量。整个前端集中于单一 app.py 文件，涵盖所有页面与交互逻辑。对于小型团队而言，该模式可有效降低维护成本。

双模融合思路的引入，源于对纯 AI 方案稳定性的测试验证。测试表明，纯 AI 方案在财务数据这种对准确性要求较高的场景下稳定性不足，需要规则进行兜底。传统指标与 AI 文本分析的组合，既保证了可解释性，又发挥了 AI 在语义理解方面的优势。

当前版本仍存在若干待改进之处：报告生成的 PDF 排版与专业排版软件产出尚存在差距；票据识别功能目前主要支持增值税发票，其他类型单据的识别效果尚不稳定。上述问题计划在后续版本中逐步优化。

综上所述，该系统已能够独立完成从数据上传、风险分析、结果展示到报告导出的完整流程。部署于云服务器后可通过公网 IP 直接访问，无需额外域名或备案。对于中小型审计机构及财务研究人员，作为辅助分析工具具有一定实用价值。

\vspace{2cm}
\begin{center}
\rule{0.6\textwidth}{0.4pt}\\[0.5em]
{\small 本文档为审计智鉴流式数据处理系统 V1.0 配套操作说明书}\\
{\small 编制日期：2026年01月}
\end{center}

\end{document}
